\documentclass[12pt,twoside]{article}
%\usepackage[table]{xcolor} % important to avoid options clash.
%\input{02465shared_preamble}
%\usepackage{cleveref}
\usepackage{url}
\usepackage{graphics}
\usepackage{multicol}
\usepackage{rotate}
\usepackage{rotating}
\usepackage{booktabs}
\usepackage{hyperref}
\usepackage{pifont}
\usepackage{latexsym}
\usepackage[english]{babel}
\usepackage{epstopdf}
\usepackage{etoolbox}
\usepackage{amsmath}
\usepackage{amssymb}
\usepackage{multirow,epstopdf}
\usepackage{fancyhdr}
\usepackage{booktabs}
\usepackage{xcolor}
\newcommand\redt[1]{ {\textcolor[rgb]{0.60, 0.00, 0.00}{\textbf{ #1} } } }


\newcommand{\m}[1]{\boldsymbol{ #1}}
\newcommand{\EE}{\mathbb{E}}
\DeclareMathOperator*{\argmax}{arg\,max}
\DeclareMathOperator*{\argmin}{arg\,min}
\newcommand{\yoursolution}{ \redt{(your solution here) } } 



\title{ Report 2 hand-in }
\date{ \today }
\author{Alice (\texttt{s000001})\and  Bob (\texttt{s000002})\and Clara (\texttt{s000003}) } 

\begin{document}
\maketitle

\begin{table}[ht!]
\caption{Attribution table. Feel free to add/remove rows and columns}
\begin{tabular}{llll}
\toprule
                                                                    & Alice   & Bob    & Clara   \\
\midrule
 1: Formulate Yodas pendulum as a linear problem                    & 0-100\%  & 0-100\% & 0-100\%  \\
 2: State at a later time                                           & 0-100\%  & 0-100\% & 0-100\%  \\
 3: State at a later time II                                        & 0-100\%  & 0-100\% & 0-100\%  \\
 4: Eigenvalues and powers                                          & 0-100\%  & 0-100\% & 0-100\%  \\
 5: Analytical expression of Eigenvalues using Euler discretization & 0-100\%  & 0-100\% & 0-100\%  \\
 6: Bound using Euler discretization                                & 0-100\%  & 0-100\% & 0-100\%  \\
 7: Matrix norm of Exponential discretization (harder)              & 0-100\%  & 0-100\% & 0-100\%  \\
 8: Stability                                                       & 0-100\%  & 0-100\% & 0-100\%  \\
 9: Discretization                                                  & 0-100\%  & 0-100\% & 0-100\%  \\
 10: Linearization                                                  & 0-100\%  & 0-100\% & 0-100\%  \\
 11: Unitgrade self-check                                           & 0-100\%  & 0-100\% & 0-100\%  \\
 12: Control using simple linearization                             & 0-100\%  & 0-100\% & 0-100\%  \\
 13: MPC                                                            & 0-100\%  & 0-100\% & 0-100\%  \\
\bottomrule
\end{tabular}
\end{table}

%\paragraph{Statement about collaboration:}
%Please edit this section to reflect how you have used external resources. The following statement will in most cases suffice: 
%\emph{The code in the irls/project1 directory is entirely}

%\paragraph{Main report:}
Headings have been inserted in the document for readability. You only have to edit the part which says \yoursolution. 

\section{Master Yodas pendulum (\texttt{yoda.py})}\label{yoda1}
\subsubsection*{{\color{red}Problem 1:  Formulate Yodas pendulum as a linear problem}}
	
	\begin{align}
		A & = \begin{bmatrix} 0 & 1 \\ -\frac{g}{L} & 0 \end{bmatrix}  \\
		B & = \begin{bmatrix} 0 \\ \frac{1}{mL^2}\end{bmatrix}
	\end{align}
	
\subsubsection*{{\color{red}Problem 2:  State at a later time}}

To solve the first part, we begin with the discretized system where $B = \mathbf{u} = 0$,
so $\dot{x} = Ax$. For Euler discretization, we approximate $\dot{x} \approx
\frac{x_{n+1} - x_n}{\Delta}$, which gives:
\begin{align}
    \frac{x_{n+1} - x_n}{\Delta} = Ax_n \implies x_{n+1} = (I + \Delta A)x_n = \tilde{A}_0 x_n
\end{align}
Applying this recursively from $x_0$:
\begin{align}
    x_1 &= \tilde{A}_0 x_0 \notag \\
    x_2 &= \tilde{A}_0 x_1 = \tilde{A}_0^2 x_0 \notag \\
    &\vdots \notag \\
    x_N &= \tilde{A}_0^N x_0
\end{align}
For exponential discretization, the exact solution to $\dot{x} = Ax$ is
$x(t) = e^{At}x_0$. Evaluating at $t = \Delta$ gives $x_1 = e^{A\Delta}x_0 = A_0 x_0$.
Applying this recursively:
\begin{align}
    x_1 &= A_0 x_0 \notag \\
    x_2 &= A_0 x_1 = A_0^2 x_0 \notag \\
    &\vdots \notag \\
    x_N &= A_0^N x_0
\end{align}
We can therefore write $x_N = \begin{bmatrix} \tilde{A}_0^N \end{bmatrix} x_0$ and
$x_N = \begin{bmatrix} A_0^N \end{bmatrix} x_0$ for Euler and exponential discretization
respectively.

As for the second part we get:
\begin{align}
\tilde A_0 & = I + \Delta A = \begin{bmatrix} 1 & \Delta \\ -\dfrac{g}{L}\Delta & 1 \end{bmatrix}, \quad A_0 = e^{A\Delta} = \text{expm}\!\left(\begin{bmatrix} 0 & \Delta \\ -\dfrac{g}{L}\Delta & 0 \end{bmatrix}\right)
\end{align}
where $\tilde{A}_0$ follows directly from $I + \Delta A$ with $A = \begin{bmatrix} 0 & 1
\\ -g/L & 0 \end{bmatrix}$, and $A_0$ is computed using the matrix exponential.

\subsubsection*{{\color{red}Problem 4: Eigenvalues and powers}}

Assume $\lambda_1, \lambda_2$ are the eigenvalues of $\tilde{A}_0$, then by the
Eigendecomposition we can write:
\begin{align}
    \tilde{A}_0 = Q\Sigma Q^{-1}, \quad \text{where } \Sigma = \begin{bmatrix} \lambda_1 & 0 \\ 0 & \lambda_2 \end{bmatrix}
\end{align}
We now consider $M = \tilde{A}_0^N$. Using the Eigendecomposition, we expand the power:
\begin{align}
    M = \tilde{A}_0^N = (Q\Sigma Q^{-1})^N
\end{align}
Writing out the multiplication explicitly:
\begin{align}
    (Q\Sigma Q^{-1})^N &= (Q\Sigma Q^{-1})(Q\Sigma Q^{-1})\cdots(Q\Sigma Q^{-1}) \notag \\
    &= Q\Sigma (Q^{-1}Q)\Sigma (Q^{-1}Q)\cdots\Sigma Q^{-1} \notag \\
    &= Q\Sigma^N Q^{-1}
\end{align}
where we used that $Q^{-1}Q = I$, which cancels all interior terms. Since $\Sigma$ is
diagonal, its $N$-th power is simply:
\begin{align}
    \Sigma^N = \begin{bmatrix} \lambda_1 & 0 \\ 0 & \lambda_2 \end{bmatrix}^N = \begin{bmatrix} \lambda_1^N & 0 \\ 0 & \lambda_2^N \end{bmatrix}
\end{align}
Therefore:
\begin{align}
    M = \tilde{A}_0^N = Q \begin{bmatrix} \lambda_1^N & 0 \\ 0 & \lambda_2^N \end{bmatrix} Q^{-1}
\end{align}
Since $M = Q\Sigma^N Q^{-1}$ is itself an Eigendecomposition, the eigenvalues of $M$
are the diagonal entries of $\Sigma^N$, namely $\lambda_1^N$ and $\lambda_2^N$.
The result is therefore very simple: \textbf{if $\lambda_1, \lambda_2$ are the
eigenvalues of $\tilde{A}_0$, then the eigenvalues of $M = \tilde{A}_0^N$ are
$\lambda_1^N$ and $\lambda_2^N$}.
	
\subsubsection*{{\color{red}Problem 5: Analytical expression of Eigenvalues using Euler discretization}}

To find the eigenvalues of $\tilde{A}_0$, we solve $\det(\tilde{A}_0 - \lambda I) = 0$.
Recall that:
\begin{align}
    \tilde{A}_0 = \begin{bmatrix} 1 & \Delta \\ -\dfrac{g}{L}\Delta & 1 \end{bmatrix}
\end{align}
We get a characteristic polynomial of:
\begin{align}
    \det(\tilde{A}_0 - \lambda I) &= \det\begin{bmatrix} 1-\lambda & \Delta \\ -\dfrac{g}{L}\Delta & 1-\lambda \end{bmatrix} \notag \\
    &= (1-\lambda)^2 + \frac{g}{L}\Delta^2 = 0
\end{align}
Rearranging:
\begin{align}
    (1-\lambda)^2 = -\frac{g}{L}\Delta^2
\end{align}
and therefore it follows from Mat1 that the two eigenvalues are:
\begin{align}
    \lambda_{1,2} = 1 \pm \Delta\sqrt{-\frac{g}{L}} = 1 \pm i\Delta\sqrt{\frac{g}{L}}
\end{align}
The eigenvalues are thus a complex conjugate pair.

\subsubsection*{{\color{red}Problem 6: Bound using Euler discretization}}

Using Euler discretization we have $x_N = \tilde{A}_0^N x_0$. Taking the norm and
using the submultiplicative property of the matrix norm:
\begin{align}
    \|x_N\| = \|\tilde{A}_0^N x_0\| \leq \|\tilde{A}_0^N\| \|x_0\| = \|\tilde{A}_0^N\|
\end{align}
where we used $\|x_0\| = 1$. From Problem 4 we know the eigenvalues of $\tilde{A}_0^N$
are $\lambda_1^N$ and $\lambda_2^N$, so by equation (12):
\begin{align}
    \|\tilde{A}_0^N\| = \max\{|\lambda_1^N|, |\lambda_2^N|\} = \max\{|\lambda_1|^N, |\lambda_2|^N\}
\end{align}
From Problem 5 we found $\lambda_{1,2} = 1 \pm i\Delta\sqrt{g/L}$, and since they are
a complex conjugate pair $|\lambda_1| = |\lambda_2|$. The absolute value of a complex
number $a + bi$ is $\sqrt{a^2 + b^2}$, so:
\begin{align}
    |\lambda_{1,2}| = \sqrt{1 + \frac{g}{L}\Delta^2}
\end{align}
Using Euler discretization we therefore get the upper bound:
$$
    \|x_N\| \leq \left(1 + \frac{g}{L}\Delta^2\right)^{N/2}
$$
This bound is tight since both eigenvalues have equal absolute value. Note that since
$1 + \frac{g}{L}\Delta^2 > 1$ for any $\Delta > 0$, the bound grows with $N$,
confirming that Euler discretization is numerically unstable for this system.

\subsubsection*{{\color{red}Problem 7: Matrix norm of Exponential discretization (harder)}}

Following the same steps as Problem 6, we start from $x_N = A_0^N x_0$ and use
$\|x_0\| = 1$:
\begin{align}
    \|x_N\| = \|A_0^N x_0\| \leq \|A_0^N\| \|x_0\| = \|A_0^N\|
\end{align}
We need the eigenvalues of $A_0^N = (e^{A\Delta})^N$. To find the eigenvalues of
$e^{A\Delta}$, we use the Eigendecomposition $A = Q\Sigma Q^{-1}$ and the Taylor
series definition of the matrix exponential:
\begin{align}
    e^{A\Delta} = \sum_{k=0}^{\infty} \frac{(A\Delta)^k}{k!} = \sum_{k=0}^{\infty} \frac{(Q\Sigma Q^{-1}\Delta)^k}{k!} = \sum_{k=0}^{\infty} \frac{Q(\Sigma\Delta)^k Q^{-1}}{k!}
\end{align}
Factoring out $Q$ and $Q^{-1}$:
\begin{align}
    e^{A\Delta} = Q\left(\sum_{k=0}^{\infty} \frac{(\Sigma\Delta)^k}{k!}\right)Q^{-1} = Q\, e^{\Sigma\Delta}\, Q^{-1}
\end{align}
Since $\Sigma = \begin{bmatrix} \mu_1 & 0 \\ 0 & \mu_2 \end{bmatrix}$ is diagonal, its
matrix exponential is simply:
\begin{align}
    e^{\Sigma\Delta} = \begin{bmatrix} e^{\mu_1\Delta} & 0 \\ 0 & e^{\mu_2\Delta} \end{bmatrix}
\end{align}
This is itself an Eigendecomposition, so the eigenvalues of $A_0 = e^{A\Delta}$ are
$e^{\mu_1\Delta}$ and $e^{\mu_2\Delta}$, where $\mu_1, \mu_2$ are the eigenvalues of
$A$. The eigenvalues of $A = \begin{bmatrix} 0 & 1 \\ -g/L & 0 \end{bmatrix}$ follow
from $\det(A - \mu I) = 0$:
\begin{align}
    \mu^2 + \frac{g}{L} = 0 \implies \mu_{1,2} = \pm i\sqrt{\frac{g}{L}}
\end{align}
So the eigenvalues of $A_0$ are $e^{\pm i\Delta\sqrt{g/L}}$, which are complex
conjugates with absolute value:
\begin{align}
    |e^{\pm i\Delta\sqrt{g/L}}| = 1
\end{align}
since $|e^{i\phi}| = 1$ for any real $\phi$ (Euler's formula: $e^{i\phi} = \cos\phi +
i\sin\phi$). From Problem 4, the eigenvalues of $A_0^N$ are then $e^{\pm
iN\Delta\sqrt{g/L}}$, which also have absolute value 1. Therefore by equation (12):
\begin{align}
    \|A_0^N\| = \max\{|e^{iN\Delta\sqrt{g/L}}|, |e^{-iN\Delta\sqrt{g/L}}|\} = 1
\end{align}
Using exponential discretization we therefore get the tight upper bound:
$$
    \|x_N\| \leq 1
$$
This result is remarkably simple and confirms that exponential discretization is
numerically stable: the norm of the state never grows beyond its initial
value, for any $\Delta > 0$ and any $N$.

\subsubsection*{{\color{red}Problem 8: Stability}}

From Problems 6 and 7 we derived the following bounds:
\begin{align}
    \text{Euler:} \quad &\|x_N\| \leq \left(1 + \frac{g}{L}\Delta^2\right)^{N/2} \\
    \text{Exponential:} \quad &\|x_N\| \leq 1
\end{align}
For Euler discretization, since $\frac{g}{L}\Delta^2 > 0$ for any $\Delta > 0$,
we have $1 + \frac{g}{L}\Delta^2 > 1$, and therefore the bound grows without bound as
$N \to \infty$:
\begin{align}
    \left(1 + \frac{g}{L}\Delta^2\right)^{N/2} \to \infty \quad \text{as } N \to \infty
\end{align}
This means the numerical solution is not guaranteed to stay bounded, regardless of
how small the time step $\Delta$ is chosen. Euler discretization is therefore
numerically unstable for this system.

For exponential discretization, the bound $\|x_N\| \leq 1$ holds for all
$N$ and all $\Delta > 0$. Since $\|x_0\| = 1$, the norm of the state never exceeds
its initial value, meaning the solution remains bounded for all time. Exponential
discretization is therefore numerically stable.

This difference in stability has a physical interpretation: the pendulum is a
system with no energy dissipation, so the true solution should neither
grow nor decay. The exponential discretization respects this property exactly, while
the Euler discretization introduces a growth that causes the solution to
break over time.

\section{R2D2 and control (\texttt{r2d2.py})}
\subsubsection*{{\color{red}Problem 9:  Discretization}}

	$$
	\m x_{k+1} = \m f_k(\m x_k, \m u_k) = \begin{bmatrix} \cdots \\ \cdots \\ \cdots \end{bmatrix}$$

\subsubsection*{{\color{red}Problem 10:  Linearization}}
	
$$
		\m x_{k+1} \approx \begin{bmatrix} \cdots \\ \cdots \\ \cdots \end{bmatrix} \m x_k + 
		\begin{bmatrix} \cdots \\ \cdots \\ \cdots \end{bmatrix} \m u_k +  
		\begin{bmatrix} \vdots \end{bmatrix} 
$$
	
\subsubsection*{{\color{red}Problem 12:  Control using simple linearization}}
			
		% Just generate the figures using the script and change the path below. 
		\begin{center}\includegraphics[width=.5\linewidth]{figures/your_answer}~
		\includegraphics[width=.5\linewidth]{figures/your_answer} \end{center}
Intuitively, the second case fails because... \yoursolution	
	
\subsubsection*{{\color{red}Problem 13:  MPC}}
			
		\begin{center}\includegraphics[width=.6\linewidth]{figures/your_answer}%~
		%	\includegraphics[width=.5\linewidth]{figures/your_answer} 
	\end{center}
		Iterative linearization solves the problem because... \yoursolution
	
\end{document}